%  #region preâmbulo
\documentclass[12pt, a4paper]{report}
\usepackage[top=3cm,left=3cm,right=2cm,bottom=2cm]{geometry}
\linespread{1.3}
\setlength{\parindent}{1.25cm}
\usepackage{indentfirst}
\usepackage[utf8]{inputenc}
\usepackage[brazil]{babel}
\usepackage{amsmath}
\usepackage{amsthm}
\usepackage{amsfonts}
\usepackage{amssymb}
\usepackage{graphicx}
\usepackage{color}
\usepackage{multicol}
\usepackage[normalem]{ulem}
\usepackage{wrapfig}
\usepackage{caption}
\usepackage{fancybox}
\usepackage[pdfstartview=FitH]{hyperref}
\usepackage{subfigure}
\bibliographystyle{plain}

\graphicspath{{Figuras/}}

\renewcommand{\theenumii}{\alph{enumii}}
\DeclareMathOperator{\sen}{sen}
\DeclareMathOperator{\tg}{tg}
\DeclareMathOperator{\arctg}{arctg}
\DeclareMathOperator{\cotg}{cotg}
\DeclareMathOperator{\agm}{agm}

\newtheorem{thm}{Teorema}[section]
\newtheorem{dfn}{Definição}[section]
\newtheorem{prob}{Problema}[section]
\newtheorem{cor}{Corolário}[section]
\newtheorem{prop}{Proposição}[section]
\newtheorem{lem}{Lema} [section]

\newcounter{contar}
%  #endregion preâmbulo

% Variáveis 
\newcommand{\nomeUniversidade}{Universidade Federal da Bahia}
\newcommand{\nomeInstituto}{Instituto de Computação}
\newcommand{\nomeCurso}{MATA53 - Teoria dos grafos}
\newcommand{\nomeProfessor}{Islame Felipe da Costa Fernandes}
\newcommand{\nomeGrupo}{\sc{\large{Antoniel Magalhães}} \\
\sc{\large{João Leahy}} \\
\sc{\large{Luis Felipe}}}
\newcommand{\titulo}{\sc{\Large{Relatório: Problema do caixeiro viajante com bônus e passageiros}}}

\begin{document}

% #region capa
\pagestyle{empty}
\begin{center}
\includegraphics[height=2.5cm]{UFBA.jpg}
\hspace{2cm}
\end{center}

\begin{center}
\sc{\large{\nomeUniversidade}} \\
\sc{\large{\nomeInstituto}} \\
\sc{\small{\nomeCurso}} \\

\vspace{4cm}

\titulo

\vspace{4.5cm}

\nomeGrupo


\vspace{5.5cm}

\textbf{Salvador - Bahia} \\
\today
\end{center}
% #endregion capa

% #region folha de rosto
\newpage
\begin{center}
\titulo

\vspace{4cm}

\nomeGrupo
\end{center}

\vspace{4cm}

\begin{flushright}
\begin{minipage}{8.6cm}
Estudo dirigido entregue ao professor \nomeProfessor
como método avaliativo da disciplina \nomeCurso


\end{minipage}
\end{flushright}
 
\vspace{8cm}


\begin{center}
\textbf{Salvador - Bahia} \\
\today
\end{center}

% #endregion folha de rosto

% #region Índice
\newpage
\tableofcontents
\thispagestyle{empty}
% #endregion Índice


% Capítulo 1: Introdução
\chapter{Introdução}

O **Problema do Caixeiro Viajante com Coleta de Bônus e Passageiros (TSP-OBP)** é uma extensão do clássico Problema do Caixeiro Viajante (TSP). Ele combina aspectos de coleta de prêmios, transporte de passageiros e otimização de rotas, tornando-se relevante em diversas aplicações, como logística, turismo e transporte público.

Este relatório tem como objetivo apresentar uma análise detalhada do problema, sua modelagem matemática, métodos de resolução e aplicações práticas.

% Capítulo 2: Fundamentação Teórica
\chapter{Fundamentação Teórica}

\section{Problema do Caixeiro Viajante (TSP)}
O Problema do Caixeiro Viajante consiste em encontrar a menor rota que visita um conjunto de cidades exatamente uma vez e retorna ao ponto de origem. É um problema clássico de otimização combinatória, amplamente estudado na teoria da computação.

\section{Extensões do TSP}
O TSP-OBP é uma extensão do TSP, incorporando os seguintes elementos:
\begin{itemize}
    \item \textbf{Coleta de Bônus}: Cada local visitado oferece um prêmio opcional. O objetivo é maximizar o bônus coletado, respeitando restrições de tempo ou custo.
    \item \textbf{Passageiros}: O transporte de passageiros pode ser incluído na rota, com cada passageiro contribuindo com uma recompensa adicional.
\end{itemize}

\section{Definição Formal do Problema}
Dado um conjunto de nós (cidades) $N$, uma matriz de custos $C = [c_{ij}]$, e bônus $b_i$ associados a cada nó $i$, o objetivo é determinar uma rota $R$ que:
\begin{itemize}
    \item Minimiza o custo total de viagem $C(R)$;
    \item Maximiza a soma dos bônus coletados $B(R)$;
    \item Satisfaz restrições, como limite de tempo e capacidade para passageiros.
\end{itemize}

% Capítulo 3: Modelagem Matemática
\chapter{Modelagem Matemática}

\section{Variáveis de Decisão}
\begin{itemize}
    \item $x_{ij} \in \{0, 1\}$: Indica se a aresta $(i, j)$ é percorrida.
    \item $y_i \in \{0, 1\}$: Indica se o bônus do nó $i$ foi coletado.
    \item $p_i \in \{0, 1\}$: Indica se um passageiro foi embarcado no nó $i$.
\end{itemize}

\section{Função Objetivo}
\[
\text{Maximizar: } \sum_{i \in N} b_i y_i - \sum_{(i,j) \in E} c_{ij} x_{ij}
\]
\section{Restrições}
\begin{enumerate}
    \item Cada nó deve ser visitado exatamente uma vez:
    \[
    \sum_{j \in N} x_{ij} = 1, \quad \forall i \in N
    \]
    \[
    \sum_{i \in N} x_{ij} = 1, \quad \forall j \in N
    \]
    \item Limites de tempo:
    \[
    \sum_{(i,j) \in E} c_{ij} x_{ij} \leq T_{\text{máximo}}
    \]
    \item Capacidade de transporte:
    \[
    \sum_{i \in N} p_i \leq \text{Capacidade}
    \]
\end{enumerate}

% Capítulo 4: Métodos de Resolução
\chapter{Métodos de Resolução}

\section{Métodos Exatos}
\begin{itemize}
    \item \textbf{Branch and Bound}: Explora todas as soluções possíveis de forma sistemática.
    \item \textbf{Programação Linear Inteira (ILP)}: Resolve a formulação matemática do problema usando técnicas de otimização.
\end{itemize}

\section{Heurísticas}
\begin{itemize}
    \item \textbf{Algoritmos Gulosos}: Constroem uma solução passo a passo, escolhendo a melhor opção local em cada etapa.
    \item \textbf{GRASP (Greedy Randomized Adaptive Search Procedure)}: Combina buscas gulosas com elementos aleatórios.
\end{itemize}

\section{Metaheurísticas}
\begin{itemize}
    \item \textbf{Algoritmos Genéticos}: Inspirados na evolução natural.
    \item \textbf{Simulated Annealing}: Baseado no recozimento térmico.
    \item \textbf{Ant Colony Optimization}: Inspirado no comportamento de colônias de formigas.
\end{itemize}

% Capítulo 5: Resultados e Discussão
\chapter{Resultados e Discussão}

\section{Resultados Computacionais}
Apresente os resultados obtidos ao resolver instâncias do problema, comparando diferentes métodos em termos de:
\begin{itemize}
    \item Tempo de execução;
    \item Qualidade da solução;
    \item Eficiência computacional.
\end{itemize}

\section{Discussão}
Discuta os trade-offs entre os métodos utilizados e os impactos práticos dos resultados.

% Capítulo 6: Conclusão
\chapter{Conclusão}
Resumo dos principais resultados e contribuições do trabalho. Apresente sugestões para trabalhos futuros, como:
\begin{itemize}
    \item Estudo de novas heurísticas para grandes instâncias.
    \item Extensão do modelo para múltiplos veículos.
    \item Aplicações em cenários reais, como planejamento logístico.
\end{itemize}

Neste trabalho visitamos teorias que fundamentam a base da computação, passando primeiro pelo concceito de funções recursivas sua importância e definições. posteriormente discorremos sobre o Lambda-cálculo e sua contextualização, e por fim relacionamos tudo apresentado com Máquinas de Turing que foi o objeto de estudo principal da disciplina.\\
É importante perceber que embora as Máquinas de Turing seja o modelo de computação preferido para o estudo em computabilidade, existem modelos clássicos dos quais o estudo continua relevante e necessário para o avanço da computação como ciência.

%-------------Bibliografia------------------
\newpage
\renewcommand{\refname}{Referências Bibliográficas}
\addcontentsline{toc}{chapter}{Referências Bibliográficas}
\bibliography{Bibliografia}
\nocite{gomes2016, gomes2018, salvador, machado, ramos, carnielli}


\end{document}
